\section{The rigidity obstruction}\label{sec:rigidity}

We now exclude the derivation supplied by \Cref{prop:extracted-lnd}.
Recall that
\[
 \Rt=\kk[a,d,b,c,u]/(ac-bd-1),\qquad
 v=a^3b-a^2u^3-d^2,
\]
and that the fiber weights of $(a,d,b,c,u)$ are
$(1,1,-1,-1,0)$.  A derivation has fiber shift $e$ if it sends
each fiber-weight-$i$ component into the fiber-weight-$(i+e)$
component.

The proof first passes to the highest part for a third, auxiliary
grading. It then studies the polynomial orbit of
$N=d^2+a^2u^3$. The polynomial $abc$ inequality forces this highest
derivation to fix $u$, and a nonsquare in the original fraction field
forces it to fix $a,d$ as well. The determinant relation leaves only derivations
whose fiber weight is positive.

The use of polynomial $abc$ to constrain locally nilpotent
derivations has precedents in
Finston--Maubach~\cite[Section~2]{FinstonMaubach}. We include the
Wronskian proof of the Mason--Stothers inequality
\cite{Mason,Stothers}, followed by the elementary one-variable fact
needed at the end of the argument.

\begin{lemma}[Mason--Stothers polynomial \(abc\)]\label{rig:abc}
Let $F$ be a field of characteristic zero, and let
$X,Y,Z\in F[q]$ be nonzero, pairwise relatively prime polynomials
satisfying $X+Y=Z$.  If they are not all constant, then
\[
 \max\{\deg X,\deg Y,\deg Z\}
 \le r(XYZ)-1,
\]
where $r(XYZ)$ is the number of distinct roots of $XYZ$ in an
algebraic closure of $F$.
\end{lemma}

\begin{proof}
Primes in this proof denote differentiation with respect to $q$.
Consider the Wronskian
\[
 W=X'Y-XY'.
\]
It is nonzero.  Indeed, if $W=0$, relative primality implies that
$X$ divides $X'$.  In characteristic zero this is impossible for
nonconstant $X$, so $X$ is a nonzero constant.  Then $Y'=0$, hence
$Y$ and $Z$ are constant as well, contrary to the hypothesis.

The same Wronskian, up to sign, is obtained from either of the
other two pairs, because $Z=X+Y$.  Work over an algebraic closure
only to count roots.  A root of multiplicity $h$ of one of the
three polynomials is not a root of either of the others.  Computing
$W$ from a pair containing that polynomial shows that $W$ vanishes
at the root to order at least $h-1$.  These roots are distinct
across the three polynomials, so, on writing
$T=\deg X+\deg Y+\deg Z$, we obtain
\[
 \deg W\ge T-r(XYZ).
\]
On the other hand, compute $W$ from the two polynomials other
than one of maximum degree $M$.  Differentiation lowers the degree
of every nonconstant polynomial by one, and therefore
\[
 \deg W\le T-M-1.
\]
Together these inequalities give $M\le r(XYZ)-1$.

The nonconstant hypothesis is necessary: for three nonzero
constants the right side would be $-1$.  The nonzero hypothesis
also causes no missing nonconstant case: if one member of a
pairwise relatively prime triple $X+Y=Z$ is zero, the other two
must be nonzero constants.  The all-zero triple is not pairwise
relatively prime.
\end{proof}

\begin{lemma}[Derivations in one variable]\label{rig:one-variable}
Let $F$ be a field of characteristic zero.  If $\delta$ is a locally
nilpotent $F$-derivation of $F[t]$, then $\delta(t)\in F$.
\end{lemma}

\begin{proof}
Write $\delta=P(t)\,d/dt$.  There is nothing to prove if $P=0$.
Otherwise let $m=\deg P$ and let $\gamma\ne0$ be its leading
coefficient.  If $m\ge1$, induction shows that, for every $n\ge1$,
the polynomial $\delta^n(t)$ has degree $1+n(m-1)$ and leading
coefficient
\[
 \gamma^n\prod_{j=0}^{n-1}\bigl(1+j(m-1)\bigr).
\]
Every factor is nonzero in characteristic zero, including when
$m=1$.  Thus no iterate vanishes, contradicting local nilpotence.
Consequently $m=0$.
\end{proof}

\begin{proposition}\label{prop:rigidity}
There is no nonzero locally nilpotent $\kk$-derivation $E$ of
$\Rt$ with a fixed fiber shift $e\le0$ and $E^2(v)=0$.
\end{proposition}

\begin{proof}
Suppose such an $E$ exists.  We introduce an \emph{auxiliary
grading}, distinct from both the valuation degree and the fiber
weight:
\begin{equation}\label{rig:auxiliary-grading}
 \begin{array}{c|rrrrr}
  &a&d&b&c&u\\ \hline
  \text{valuation degree}&0&0&0&0&0\\
  \text{fiber weight}&1&1&-1&-1&0\\
  \text{auxiliary degree}&0&3&-3&0&2.
 \end{array}
\end{equation}
The relation $ac-bd-1$ is homogeneous of auxiliary degree zero,
so this is an actual $\mathbb Z$-grading on $\Rt$.  Its components
need not vanish in negative degrees.  The generators are
homogeneous for both this grading and the fiber grading, and the
relation is homogeneous for both; thus the two gradings commute.

\medskip\noindent
\emph{The highest auxiliary component.}
For an auxiliary-homogeneous $r\in\Rt$ of degree $i$, define
$E_j(r)$ to be the degree-$(i+j)$ component of $E(r)$, and extend
linearly to $\Rt$.  Projecting the product rule to each degree
shows that every $E_j$ is a derivation.  Only finitely many of
these derivations are nonzero: their values on the five generators
are obtained from the finitely many homogeneous components of
their images under $E$, and a derivation vanishing on all five
generators is zero.  We therefore have a finite decomposition
\[
 E=\sum_j E_j.
\]
Let $\lambda$ be its greatest nonzero shift and set $E'=E_\lambda$.
Then $E'\ne0$.  For every auxiliary-homogeneous $r$ of degree $i$
and every $n\ge1$, the component of $E^n(r)$ of degree
$i+n\lambda$ equals $(E')^n(r)$: in the expansion of the iterate,
each chosen shift is at most $\lambda$, and their sum can equal
$n\lambda$ only when every chosen shift equals $\lambda$.
Local nilpotence of $E$ therefore implies local nilpotence of
$E'$ on every homogeneous $r$, and hence on every element, which
is a finite sum of such elements.  This reasoning uses finite
shift sets and finite sums, so the negative degrees in
\Cref{rig:auxiliary-grading} cause no difficulty.  Projection
to auxiliary degree preserves fiber weight, and consequently
$E'$ still has fiber shift $e$.

Set
\[
 N=d^2+a^2u^3.
\]
Both its summands have auxiliary degree $6$, whereas $a^3b$ has
auxiliary degree $-3$.  In $E^2(v)=E^2(a^3b-N)$, the first term
has auxiliary degrees at most $-3+2\lambda$.  The degree
$6+2\lambda$ component of the second term is $(E')^2(N)$.
Taking that component of $E^2(v)=0$ gives
\begin{equation}\label{rig:linear-orbit}
 (E')^2(N)=0.
\end{equation}

\medskip\noindent
\emph{The orbit polynomials fix $a,d,u$.}
The localization obtained by solving for $c$ is
\begin{equation}\label{rig:a-chart}
 \Rt[a^{-1}]=\kk[a,a^{-1},d,u][b],
 \qquad c=\frac{1+bd}{a}.
\end{equation}
In particular, $a,d,b,u$ are algebraically independent,
$N\ne0$, and the fraction field is
\begin{equation}\label{rig:original-field}
 K=\Frac(\Rt)=\kk(a,d,b)(u).
\end{equation}
Since $E'$ is locally nilpotent, exponentiating it gives
polynomials
\[
 a(q)=\exp(qE')(a),\quad
 d(q)=\exp(qE')(d),\quad
 u(q)=\exp(qE')(u)
 \quad\text{in }\Rt[q]\subset K[q].
\]
Each is nonzero, since its value at $q=0$ is the corresponding
nonzero element of $\Rt$.  The exponential is multiplicative
by the iterated product rule.  Hence \Cref{rig:linear-orbit}
implies
\begin{equation}\label{rig:orbit-sum}
 d(q)^2+a(q)^2u(q)^3
 =N(q):=\exp(qE')(N)=N+qE'(N),
 \qquad \deg_q N(q)\le1.
\end{equation}
Here $N(q)\ne0$ because $N(0)=N\ne0$.  Thus neither summand nor
their sum is the zero polynomial.

The three polynomials in \Cref{rig:orbit-sum} are pairwise
relatively prime in $K[q]$.  To see this, a common root of the
two summands over $\overline K$ would have multiplicity at least
two in each summand.  For the second summand this follows
whether the root belongs to $a(q)$, to $u(q)$, or to both.
Their sum would then have a root of multiplicity at least two,
which a nonzero polynomial of degree at most one cannot have.
Thus the summands are relatively prime, and the identity
$X+Y=Z$ implies relative primality of the other two pairs.
No relative primality between $a(q)$ and $u(q)$ is needed.

Write
\[
 D_a=\deg_q a(q),\qquad D_d=\deg_q d(q),\qquad
 D_u=\deg_q u(q).
\]
Suppose at least one summand in \Cref{rig:orbit-sum} is
nonconstant.  A nonconstant summand has degree at least two.
Since the sum has degree at most one, both summands must then
have the same degree
\[
 M=2D_d=2D_a+3D_u\ge2.
\]
The number of distinct roots of their product with $N(q)$ is
at most $D_d+D_a+D_u+\deg_q N(q)$.  Applying \Cref{rig:abc}
therefore yields
\begin{equation}\label{rig:degree-bound}
 2D_d=2D_a+3D_u
 \le D_d+D_a+D_u+\deg_q N(q)-1
 \le D_d+D_a+D_u.
\end{equation}
This also covers constant $N(q)$; in that case the first bound
is stronger.  The last inequality gives $D_d\le D_a+D_u$.
Combining it with $2D_d=2D_a+3D_u$ gives $3D_u\le2D_u$,
so $D_u=0$.  Consequently $u(q)=u$ and
$D_d=D_a>0$.

Let $\alpha,\delta\in K^*$ be the leading coefficients of
$a(q)$ and $d(q)$, respectively.  Cancellation of the
degree-$M$ terms in \Cref{rig:orbit-sum} requires
\[
 \delta^2+\alpha^2u^3=0,
 \qquad\text{hence}\qquad
 \left(\frac{\delta}{\alpha}\right)^2=-u^3.
\]
This equality would hold in the original field $K$ of
\Cref{rig:original-field}, not merely in $\overline K$.
It is impossible: put $K_0=\kk(a,d,b)$, so that $K=K_0(u)$.
For a nonzero rational function of $u$, the multiplicity of $u$
in its numerator minus its multiplicity in its denominator
defines an integer-valued order.  A square has even order,
while $-u^3$ has order three.  Passing to $\overline K$ above
was used only for counting roots and has not changed the field
of the leading coefficients.

We conclude that both summands in \Cref{rig:orbit-sum} are
constant.  Their nonvanishing gives $2D_d=0$ and
$2D_a+3D_u=0$.  All three degrees are nonnegative, so
$D_a=D_d=D_u=0$.  Taking coefficients of $q$ in their
exponentials now gives
\begin{equation}\label{rig:fixed-coefficients}
 E'(a)=E'(d)=E'(u)=0.
\end{equation}

\medskip\noindent
\emph{The remaining coefficient has an impossible fiber weight.}
Differentiating $ac-bd=1$ and using
\Cref{rig:fixed-coefficients} gives $aE'(c)=dE'(b)$.  Define
\[
 h=cE'(b)-bE'(c)\in\Rt.
\]
Using both determinant identities, without dividing by either
$a$ or $d$, gives
\begin{equation}\label{rig:h-formulas}
 a h=(ac-bd)E'(b)=E'(b),\qquad
 d h=(ac-bd)E'(c)=E'(c).
\end{equation}
Let $C_0=\kk[a,d,u]$ and $Q=\Frac(C_0)=\kk(a,d,u)$.
By \Cref{rig:fixed-coefficients}, every nonzero element of
$C_0$ is invariant.  Thus $E'$ extends to a locally nilpotent
$Q$-derivation of
\[
 (C_0\setminus\{0\})^{-1}\Rt=Q[b],
\]
where the equality follows by solving for $c$ as in
\Cref{rig:a-chart}.  Local nilpotence is preserved because
$(E')^n(r/s)=(E')^n(r)/s$ for $s\in C_0\setminus\{0\}$.
By \Cref{rig:one-variable}, $E'(b)\in Q$; hence
\Cref{rig:h-formulas} implies $h\in Q$.

We spell out the coefficient intersection that improves this
to $h\in C_0$.  Besides \Cref{rig:a-chart}, solving for $b$
gives
\[
 \Rt[d^{-1}]=\kk[a,d,d^{-1},u][c],
 \qquad b=\frac{ac-1}{d}.
\]
Since $b$ and $c$ are each transcendental over $Q$, comparison
of polynomial coefficients in these two presentations gives
\[
 \Rt\cap Q
 \subset C_0[a^{-1}]\cap C_0[d^{-1}]
 =C_0.
\]
The last intersection is taken in $Q$: in the unique
factorization domain $C_0$, the reduced denominator of an
element belonging to both rings must divide a power of $a$
and a power of $d$.  As $a$ and $d$ are relatively prime, that
denominator is a unit.  Conversely $C_0\subset\Rt\cap Q$, so
the intersection is exactly $C_0$.  In particular $h\in C_0$.

The derivation $E'$ is nonzero and fixes $a,d,u$; therefore
\Cref{rig:h-formulas} implies $h\ne0$.  Since $E'$ has fiber
shift $e$, the definition of $h$ shows that it is homogeneous
of fiber weight $e-2<0$.  But every monomial in
$C_0=\kk[a,d,u]$ has nonnegative fiber weight, because the
weights of $a,d,u$ are $1,1,0$.  A nonzero element of $C_0$
cannot have negative fiber weight.  This contradiction proves
the proposition.
\end{proof}

\begin{proof}[Proof of \Cref{thm:main}]
The explicitly defined algebra $A=P/(H)$ is a finitely
generated unital $\kk$-algebra.  By \Cref{con:domain} it is
integral, with fraction field $\kk(p,x,y,z)$, and hence has
dimension four.  Moreover, \Cref{prop:stabilization}
gives a $\kk$-algebra isomorphism
\[
 A[w]\cong\kk\poly{5}.
\]
Suppose that $A\cong\kk\poly{4}$.  By
\Cref{prop:graded,prop:positive-invariant}, its associated
graded ring $G$ admits a nonzero homogeneous locally nilpotent
derivation $D_0$ with a nonzero homogeneous invariant of
positive valuation degree.  The principal bundle and its
coordinate ring $\Gt$ are given by
\Cref{prop:principalization}.  By
\Cref{prop:homogeneous-lift}, $D_0$ lifts to a locally
nilpotent derivation $\Dt$ on $\Gt$ with the same valuation
degree shift and fiber shift zero.  The reduction of this lift
and invariant in \Cref{prop:extracted-lnd} produces
a nonzero locally nilpotent derivation $E$ of $\Rt$ with a
fixed nonpositive fiber shift and $E^2(v)=0$.  This contradicts
\Cref{prop:rigidity}.  Thus $A\not\cong\kk\poly{4}$, while
its cylinder is $\kk\poly{5}$.  Taking $n=4$ proves the
claimed counterexample over $\kk$.
\end{proof}
